\documentclass[11pt,a4paper]{article}
\usepackage[utf8]{inputenc}
\usepackage[T1]{fontenc}
\usepackage[spanish]{babel}
\usepackage{lmodern}
\usepackage[a4paper,margin=23mm]{geometry}
\usepackage{xcolor,graphicx,hyperref,booktabs,longtable,array,fancyhdr,float,listings,enumitem}
\usepackage{microtype}
\definecolor{azul}{HTML}{1A4F74}
\hypersetup{colorlinks=true,linkcolor=azul,urlcolor=azul}
\setlength{\headheight}{15pt}\pagestyle{fancy}\fancyhf{}\lhead{uCuPViewer | Verificación}\rhead{Milk-V Duo S + STM32L432KC}\cfoot{\thepage}
\lstset{basicstyle=\ttfamily\small,breaklines=true,frame=single,columns=fullflexible,showstringspaces=false}
\newcommand{\Pendiente}{\textcolor{red!60!black}{\textbf{PENDIENTE}}}
\newcommand{\Registro}[3]{%
\subsection{#1 -- #2}
\noindent\textbf{Requisitos relacionados:} #3\\
\textbf{Estado:} \Pendiente\quad \textbf{Responsable:} \underline{\hspace{35mm}}\\
\textbf{Fecha / commit:} \underline{\hspace{60mm}}\\
\textbf{Instrumento o herramienta:} \underline{\hspace{75mm}}
\paragraph{Procedimiento ejecutado} Describir exactamente los pasos realizados, la configuración y número de repeticiones.
\paragraph{Resultados} Indicar valor objetivo, medida, unidad, tolerancia adoptada previamente, cálculo y dictamen.
\paragraph{Evidencia} Enlazar rutas o insertar una imagen/captura y señalar qué demuestra; no inferir más de lo medido.\par\medskip
\noindent\textbf{Archivo/URL:} \texttt{tests/evidence/\ldots}\\
\textbf{Conclusión y límites:} \underline{\hspace{95mm}}\par\bigskip}
\newcommand{\EvidenciaImagen}[2]{\begin{figure}[H]\centering\IfFileExists{#1}{\includegraphics[width=.90\linewidth]{#1}}{\fbox{\parbox[c][35mm][c]{.83\linewidth}{\centering EVIDENCIA PENDIENTE\\\texttt{\detokenize{#1}}}}}\caption{#2}\end{figure}}
\title{\textbf{Informe de verificación experimental y documental}\\\large uCuPViewer: Milk-V Duo S + STM32L432KC}
\author{Equipo: \underline{\hspace{45mm}}}
\date{Base revisada: main 41601d45f7\\ Fecha de cierre: \underline{\hspace{33mm}}\\ Versión del contrato: \underline{\hspace{24mm}}}
\begin{document}
\maketitle
\begin{abstract}
\noindent\textbf{Estado de esta plantilla:} sin resultados inventados. Completar únicamente con medidas y revisiones realizadas. Diferenciar simulación, compilación, prueba en placa y ensayo físico del sistema integrado.
\end{abstract}
\tableofcontents\newpage
\section{Identificación del montaje y trazabilidad}
\begin{longtable}{p{50mm}p{105mm}}\toprule
Campo & Valor documentado\\\midrule
SBC / revisión & Milk-V Duo S / \underline{\hspace{45mm}}\\
Linux (kernel / imagen) & \underline{\hspace{80mm}}\\
CPU/arquitectura de arranque & \underline{\hspace{80mm}}\\
MCU / placa & STM32L432KC / \underline{\hspace{45mm}}\\
Firmware (SHA y compilador) & \underline{\hspace{80mm}}\\
Backend/frontend (SHA) & \underline{\hspace{80mm}}\\
SPI: GPIO, modo, pinout & \underline{\hspace{80mm}}\\
Servidor web / versión & \underline{\hspace{80mm}}\\
Instrumentación / calibración & \underline{\hspace{80mm}}\\
\bottomrule\end{longtable}

\section{Método y criterios adoptados}
El docente exige dos canales, ventanas de al menos 64 muestras, presets 100, 1.000 y 10.000 Sa/s, muestreo constante, SPI, servidor HTTP y refresco entre 1 y 4 segundos. Para verificación conservadora se adoptan 64 muestras y las tasas por canal (confirmar con el docente).

\subsection{Estado y contrato de la implementación actual}
La implementación integrada usa 64 pares CH1/CH2 de 12 bits almacenados como
\texttt{uint16}, cabecera SC de seis bytes y CRC16 de dos bytes: 264 bytes en
cada transferencia. El primer byte del maestro solicita un preset de la próxima
captura, sin ACK ni READY independientes. El STM32 actual usa TIM2 y
\emph{polling} ADC; el código ADC+DMA previo se conserva en otra ruta.

El presupuesto anterior de 200 ms por SPI fue una \emph{propuesta de ingeniería}
y no una exigencia académica ni resultado demostrado. En una sesión física el
maestro C completó dos lecturas de 264 B en 1,783 y 1,791 s. La
frecuencia real de SCLK requiere un instrumento de medición. No confundir estas
latencias con la tasa de conversión ADC.

\[
  R_{\rm util} = \frac{256\ \mathrm{bytes}}{T_{\rm SPI}},
  \qquad
  \varepsilon_{f} = 100\frac{f_{\rm medida}-f_{\rm objetivo}}{f_{\rm objetivo}}.
\]
Las tasas obligatorias nominales 100, 1.000 y 10.000 Sa/s por canal requieren
medición externa y tolerancia acordada antes de dictamen.

\section{Pruebas físicas y funcionales}
\Registro{F-HW01}{Montaje, niveles eléctricos y asignación de pines}{RD-05 / R-05}
\Registro{F-AD01}{Dos entradas analógicas distinguibles}{RV-02}
\Registro{F-AD02}{Integridad de ventanas y longitud mínima}{RV-03 / RV-04}
\Registro{F-AD03}{Frecuencias nominales 100, 1.000 y 10.000 Sa/s}{RV-05 / RV-06 / RV-07}
\Registro{F-AD04}{Constancia de muestreo y variación temporal}{RV-08}
\Registro{F-AD05}{Desfase entre canales multiplexados}{RD-06}
\Registro{F-AD06}{Bonificación: 40.000 Sa/s}{RV-27 / BON-01}
\Registro{F-SP01}{SPI por GPIO: modo, secuencias y niveles}{RV-10 / RD-09}
\Registro{F-SP02}{Integridad de tramas SC264, CRC, secuencia y errores}{RD-03 / RD-04}
\Registro{F-SP03}{Rendimiento real y repetibilidad de SPI por GPIO}{RV-11 / RV-12 / RV-13 / RD-12}
\noindent\textbf{Tabla de repetibilidad SPI (datos medidos):}
\begin{center}\begin{tabular}{lrrrr}\toprule
Preset & N & Media [ms] & p95 [ms] & Máx. [ms]\\\midrule
100 Sa/s & -- & -- & -- & --\\
1.000 Sa/s & -- & -- & -- & --\\
10.000 Sa/s & -- & -- & -- & --\\\bottomrule\end{tabular}\end{center}
\Registro{F-SP04}{Cambio real del preset desde la interfaz SPI}{RV-09}
\Registro{F-WEB01}{Dashboard HTTP servido desde Milk-V, no solo PC}{RV-14 / RV-15}
\Registro{F-WEB02}{Cadencia de actualización gráfica}{RV-16}
\Registro{F-WEB03}{Estados visibles y datos no simulados}{RD-07}
\Registro{F-IN01}{Integración hasta HTTP en SBC y navegador final}{RV-01 / RV-20}
\Registro{F-IN02}{Robustez ante desconexión y errores}{RD-11}
\Registro{F-SYS01}{Arranque y servicio web reproducibles en Milk-V}{RD-10}

\section{Verificaciones de código, configuración y documentos}
La evidencia puede consistir en un \textbf{fragmento de código, nombre de archivo, commit y una prueba de compilación} (no es obligatorio un instrumento físico). Registrar qué demuestra y qué NO demuestra cada fragmento.

\subsection{Implementación actual: TIM2 y adquisición secuencial}
\noindent\textbf{Ruta:} \texttt{firmware/stm32/spi\_polling/src/main.c}\\
\textbf{Base:} \texttt{41601d45f7}. \textbf{Símbolos:} \texttt{rate\_hz},
\texttt{prepare\_frame}, \texttt{slave\_transfer}.\\
\textbf{Demuestra:} selección de períodos nominales y detección de periodos
perdidos; no demuestra la frecuencia efectiva.
\begin{lstlisting}[language=C]
uint32_t period_us = 1000000u / rate_hz();
htim2.Init.Prescaler = 15;
htim2.Init.Period = period_us - 1u;
/* La lectura ADC se inicia por software tras observar TIM2. */
\end{lstlisting}
\noindent\textbf{Referencia fuente:}
\href{https://github.com/anhenaobe/uCuPViewer/blob/41601d45f71fb994444670fc85fe28a7be1d0a50/firmware/stm32/spi_polling/src/main.c}{Código STM32 en GitHub, commit 41601d45f7}

\EvidenciaImagen{../evidence/docs/ejemplo_build.png}{Ejemplo de imagen de compilación: insertar evidencia real y explicar su alcance.}
\Registro{D-SPI01}{Teoría de SPI y plataforma}{RV-11/12, RD-09}
\Registro{D-SP02}{Contrato SPI aprobado}{RD-03/04}
\Registro{D-WEB01}{Servicio web implementado}{RV-15}
\Registro{D-TEST01}{Pruebas unitarias}{RV-19}
\Registro{D-DOC01}{Equipo y declaraciones}{RV-17/18/25/26}
\Registro{D-DOC02}{Especificación y reproducibilidad}{RV-21/22/23}
\Registro{D-DOC03}{Video y demostración}{RV-24}

\section{Pruebas previas localizadas en GitHub, sin repetición en este informe}
\subsection{SPI físico: validación limitada}
Registro documental del proyecto: dos transferencias del maestro C de 264 bytes,
con tiempos 1,783 y 1,791 s, secuencias 1 y 2; una transferencia Python adicional
secuencia 3. Las tres fueron registradas con cabecera SC y CRC válidos.
\url{https://github.com/anhenaobe/uCuPViewer/blob/main/docs/debug/milkv_validacion_post_systick.md}
\textbf{Límite:} no se midió aquí SCLK ni tasa ADC ni se realizó un ensayo de 100 ventanas.

\subsection{HTTP provisional en PC/VM y fallos observados}
Registro documental: 64 pares por captura, API con \texttt{source=hardware},
37 capturas válidas y dos errores de JSON UART; un error \texttt{0x17} en
el preset nominal de 10.000 Sa/s seguido de recuperación al preset 1.000.
El HTTP observado se servía desde PC/VM, no desde Milk-V con Apache/Nginx.
\url{https://github.com/anhenaobe/uCuPViewer/blob/main/docs/debug/milkv_dashboard_uart_validacion.md}

\section{Video de demostración: 100, 1000 y 10000 Sa/s}
\noindent\textbf{Enlace real a Google Drive:}\par\smallskip
\noindent\fbox{\parbox{.88\linewidth}{\centering \textbf{PENDIENTE: PEGAR AQUÍ URL DE DRIVE}\par
\vspace{5mm}\hrulefill}}\par\medskip
\noindent\textbf{Responsable / fecha / SHA:} \underline{\hspace{65mm}}\\
\textbf{Permisos comprobados por:} \underline{\hspace{65mm}}\par\medskip
\begin{center}
\begin{tabular}{p{27mm}p{26mm}p{92mm}}
\toprule
Modo nominal & Momento video & Evidencia obligatoria a mostrar\\\midrule
100 Sa/s & --:-- & 64 pares por canal, señales reales y medición temporal\\
1000 Sa/s & --:-- & 64 pares, señales reales y medición temporal\\
10000 Sa/s & --:-- & 64 pares, señales reales, manejo de 0x17 y medición temporal\\
Servidor / SPI & --:-- & Milk-V, STM32, trama y HTTP en SBC si está desplegado\\
\bottomrule
\end{tabular}
\end{center}
\textbf{Advertencia:} ver 10000 en una pantalla no demuestra que el ADC alcance
10000 muestras/s. Anexar instrumento, método y datos a \texttt{tests/evidence/adc/}.
La guía editable y casilla central del video está en
\texttt{docs/VERIFICACION\_CODIGO\_Y\_EVIDENCIAS.md}, sección 8.

\section{Consolidación y dictamen}
\begin{longtable}{p{20mm}p{24mm}p{20mm}p{72mm}}\toprule
Ensayo & Estado (P/C/NC/B) & Fecha & Evidencia / nota\\\midrule\endhead
F-HW01 & P & -- & --\\
F-AD01 & P & -- & --\\
F-AD02 & P & -- & --\\
F-AD03 & P & -- & --\\
F-AD04 & P & -- & --\\
F-AD05 & P & -- & --\\
F-AD06 & P & -- & --\\
F-SP01 & P & -- & --\\
F-SP02 & P & -- & --\\
F-SP03 & P & -- & --\\
F-SP04 & P & -- & --\\
F-WEB01 & P & -- & --\\
F-WEB02 & P & -- & --\\
F-WEB03 & P & -- & --\\
F-IN01 & P & -- & --\\
F-IN02 & P & -- & --\\
F-SYS01 & P & -- & --\\
D-SPI01 & P & -- & --\\
D-SP02 & P & -- & --\\
D-WEB01 & P & -- & --\\
D-TEST01 & P & -- & --\\
D-DOC01 & P & -- & --\\
D-DOC02 & P & -- & --\\
D-DOC03 & P & -- & --\\
\bottomrule\end{longtable}

\section{Limitaciones, incidentes y acciones correctivas}
Registrar mediciones inconclusas, condiciones de laboratorio, fallos encontrados, issue/PR correctivo, decisión de repetición y cambios de tolerancias. No eliminar las evidencias negativas.

\section{Acta de revisión}
\noindent\textbf{Responsable de prueba:} \underline{\hspace{50mm}}\\
\textbf{Revisor técnico:} \underline{\hspace{50mm}}\\
\textbf{Fecha y SHA del cierre:} \underline{\hspace{55mm}}\\
\textbf{Ubicación del video:} \underline{\hspace{75mm}}

\end{document}
